\section{The Cusp Is Physical Space}

The claim of this part is direct. Ordinary three-dimensional space is the flat $\{4,3,4\}$ cubic cusp of the
four-dimensional bulk. The bulk is the hyperbolic crystal $\{3,4,3,4\}$ of $24$-cells. At each ideal corner of
that crystal sits a Euclidean cubic honeycomb, and this section argues that one such cusp is the space we live
in. The work here is to say precisely what space is on this proposal and to show by measurement that it is
genuinely and cleanly three-dimensional. Two pictures recur, the burning bulk and its cast radiance, but the
substance is in the geometry and the numbers.

\subsection{Paracompactness is a feature}

The honeycomb $\{3,4,3,4\}$ is \emph{paracompact}. Its vertex figure is the Euclidean $\{4,3,4\}$ cubic
honeycomb, an ideal element with infinitely many cells rather than a finite polytope. In the usual classification
of regular hyperbolic honeycombs this counts as a degenerate case, since the element sits at infinity. We read it
the other way. That ideal cubic element is the flat three-dimensional physical space, and its appearing at the
ideal corner is exactly what hands physics a clean, periodic, flat crystal to live on.

Seen this way the paracompactness is not a blemish to be apologized for. It is the mechanism. A genuinely compact
honeycomb has finite cells, but the honeycombs of that kind that exist in four dimensions are all
non-crystallographic, and so they carry no spinor structure on their direction set. The compact options buy a
bounded vertex figure at the cost of the spinor coin. The paracompact $D_4$ cusp pays the opposite price and
gains the opposite prize. It surrenders a finite vertex figure and in exchange it gets both the spinor coin in
the bulk and an exactly flat cubic three-space at the corner.

\begin{proposition}[The ideal element is physical space]
The vertex figure of $\{3,4,3,4\}$ is the Euclidean cubic honeycomb $\{4,3,4\}$. Read as the geometry at the
ideal corner of the bulk, this paracompact element is the flat three-dimensional space on which matter
propagates. Paracompactness is therefore the feature that supplies physical space, not a defect of the
substrate.
\end{proposition}

\subsection{Flat at finite distance, not only at infinity}

A natural worry is that a structure living at an ideal corner is real only in a limit, never at any finite stage.
The geometry rules this out. A horosphere in hyperbolic space is intrinsically Euclidean at every finite level.
This is a theorem about hyperbolic geometry, not an asymptotic statement, and the cubic cusp inherits it. The
flatness holds on a real surface at finite distance, not only in a limit one approaches and never reaches.

The same point shows up in measurement. A horosphere band, grown outward, expands polynomially with an exponent
near $1.5$, while the surrounding bulk expands exponentially with a rate near $9$ \cite{pollard2026vibetest}. The
gap is unmistakable. The bulk swells the way hyperbolic volume swells, and the horospherical surface inside it
grows the slow way a flat three-space grows. The flat cusp is therefore a real, finite-distance, exactly
Euclidean surface sitting inside an exponentially growing curved bulk, not an idealization that exists only at
the boundary.

\begin{result}[The cusp is flat at finite distance]
The cubic cusp grows polynomially, with measured band exponent near $1.5$, against the bulk's exponential rate
near $9$. Combined with the horosphere flatness theorem, this confirms an exactly Euclidean three-dimensional
surface at every finite stage, not merely in the ideal limit \cite{pollard2026vibetest}.
\end{result}

\subsection{The clean cusp versus a generic slice}

There is a sharper question than flatness. Among the surfaces one could carve out of the bulk, which one is
physical space? Two candidates compete. One is the clean $\{4,3,4\}$ cubic cusp, the periodic crystal sitting at
the ideal corner. The other is a generic horosphere slice, an arbitrary level surface cut through the bulk with
no regard for the cubic lattice. An earlier and more optimistic reading hoped that a generic aperiodic slice
might be the better physical space, on the thought that aperiodicity buys isotropy for free. The measurement did
not support that reading.

\begin{table}[h]
\centering
\begin{tabular}{lcc}
\toprule
property & clean cubic cusp & generic horosphere slice \\
\midrule
periodicity & periodic crystal & aperiodic, rough \\
spectral dimension & $2.97$ (clean 3D) & $2.16$ (sub-3D) \\
\bottomrule
\end{tabular}
\end{table}

The clean cubic cusp reads a spectral dimension of $2.97$, which is clean three-dimensional behaviour within
measurement error. The generic slice reads $2.16$, distinctly sub-three-dimensional, the signature of a rough
aperiodic surface where a diffusing walker keeps getting trapped \cite{pollard2026vibetest}. The verdict is that
the clean cubic cusp, not the disordered generic slice, is physical space. The aperiodic-is-isotropic intuition
was attractive but wrong, and the test settled it in favour of the periodic crystal.

\begin{result}[The clean cusp is clean 3D]
The $\{4,3,4\}$ cubic cusp measures spectral dimension $2.97$, clean three-dimensional, while a generic
horosphere slice measures $2.16$, sub-three-dimensional. Physical space is the clean cubic cusp, not a generic
slice \cite{pollard2026vibetest}.
\end{result}


This pins down what space is on the proposal. It is the clean $\{4,3,4\}$ cubic cusp of the bulk, genuinely and
cleanly three-dimensional, exactly flat at finite distance, the radiance cast by the burning bulk rather than a
limiting fiction. The paracompactness that looked like a defect is the very thing that supplies a flat crystal
for matter. One residual remains. The cusp carries a small cubic anisotropy, the order-four imperfection of the
octahedral point group, which the later audit treats in full. Everything else about the identity of space is
settled here, and settled cleanly.
